% !TEX root = 00_dmlr_compel.tex
\section{Introduction}\label{sec:intro}

Recent advances in language models have been driven by pretraining on massive web-scale corpora \citep{brown2020language, du2022glam, hoffmann2022training, raffel2020exploring, xie2023dsir}.
However, the sheer size of these datasets---often comprising trillions of tokens---makes efficient data selection both crucial and challenging \citep{albalak2024survey}.

Carefully curated pretraining data improves downstream performance, as shown by corpora like FineWeb-EDU \citep{penedo2024fineweb} and DCLM \citep{li2024datacomp}.
These datasets rely on multi-stage filtering pipelines that combine language identification, duplication and toxicity checks, and scoring heuristics based on LLMs or perplexity \citep{gehman2020realtoxicity, lee2022deduplicating, marion2023less, albalak2024survey}.
While these pipelines are effective, the notion of ``high quality'' remains under-defined, and existing filters often reflect implicit or task-specific assumptions \citep{longpre2023flan}.
As pretraining datasets grow to tens of trillions of tokens, we need fast and interpretable signals to help identify high-quality data at scale.

This raises a core question: \emph{How can we identify high quality data cheaply and quickly?}

\begin{figure}[t]
\centering
\includegraphics[width=\textwidth]{fig1_pipeline.png}
\caption{\textbf{\textsc{Compel} filters low-quality documents using compression ratio.}
(1) Start with a web-scale input corpus.
(2) Compute the compression ratio (CR) of each document using LZ4 compression.
(3) Discard documents with CR outside a calibrated quality band (e.g., $[0.65, 0.80]$).
(4) Retain the filtered corpus for pretraining.
(5) Pretrain language models with the improved dataset.
\textsc{Compel} improves data quality using a fast, model-free signal that requires no training, inference, or supervision.}
\label{fig:pipeline}
\end{figure}

There are likely many promising answers to this question.
In this work, we introduce an approach that is complementary to existing methods and integrates easily into current filtering pipelines.
We propose \textsc{Compel}, a scalable and lightweight filtering method that uses a document's compression ratio as a proxy for information density (Figure~\ref{fig:pipeline}).
In information theory, data compression encodes information using fewer bits than the original representation, reducing redundancy without losing essential content \citep{shannon1948mathematical}.
Our key insight is that data compression provides a surprisingly effective signal for identifying high-quality text.
Intuitively, well-edited, information-rich documents contain meaningful content that is neither overly repetitive nor excessively noisy.
These documents tend to compress moderately.
In contrast, boilerplate or repetitive content compresses too well (e.g., thousands of identical HTML tags), while noisy or unnatural data compresses poorly due to its lack of structure.
By discarding documents with extremely low or high compression ratios, \textsc{Compel} improves the overall information density of a corpus.
We calibrate filtering thresholds using high-quality reference datasets and show that this lightweight filter can further enhance performance---even when applied on top of state-of-the-art heuristic or model-based filtering.

We validate our method on three leading pretraining corpora---FineWeb, FineWeb-EDU, and DCLM---each of which employs a distinct filtering strategy.
Across 13 downstream tasks, models trained on \textsc{Compel}-filtered data consistently outperform those trained on the original corpora.
These improvements are robust and achieved at negligible computational cost.
The resulting accuracy gains are modest (+0.4--1.1 points), but small, consistent gains from lightweight filters are substantial at scale, offering considerable efficiency benefits on trillion-token datasets \citep{raffel2020exploring, wenzek2020ccnet, ankner2024perplexed}.

\paragraph{Contributions.}
This work contributes to the data side of language-model development:
\begin{itemize}[leftmargin=1.35em,topsep=2pt,itemsep=1pt]
\item A distributional analysis of LZ4 compression ratios across four corpora of varying curation quality (C4, FineWeb-EDU, DCLM, Twitter), showing that higher-quality corpora concentrate in a narrow mid-range band and that this band mirrors the accept/reject decisions of the FineWeb-EDU quality classifier (Section~\ref{sec:cr-signal}).
\item An information-theoretic account of why an interior compression-ratio band should select learnable text, connecting \textsc{Compel} to Shannon entropy, Kolmogorov complexity, and the minimum description length principle (Section~\ref{sec:theory}).
\item \textsc{Compel}, a model-free filter that retains documents whose compression ratio falls in a calibrated band ($[0.65, 0.80]$) and processes over 40{,}000 documents per second on 500 CPUs (Sections~\ref{sec:method}).
\item Pretraining validation at 1.4B and 8B parameters on FineWeb, FineWeb-EDU, and DCLM: macro-average accuracy over 13 benchmarks improves by +1.1, +0.7, and +0.4 points respectively at 1.4B scale, and macro and micro averages also improve at 8B scale (Section~\ref{sec:results}).
\end{itemize}
